\documentclass[10pt, twoside]{article}
\usepackage[margin=1in]{geometry}
\usepackage{setspace}
\usepackage{fancyhdr}

\geometry{
    paperwidth=6in,  % Adjust the width of the paper
    paperheight=9in,   % Adjust the height of the paper
    margin=0.85in,
}

% Line spacing: tighter than double
\setstretch{1.2}

\setlength{\parindent}{1.5em}
\setlength{\parskip}{0pt}

\newcommand{\essaytitle}{Sermon on Death}
\newcommand{\essayauthor}{Jacques-Bénigne Bossuet}

\title{\essaytitle}
\author{\essayauthor}

\makeatletter
\let\thetitle\@title
\let\theauthor\@author
\makeatother

% Header style
\pagestyle{fancy}
\fancyhf{}
\fancyhead[LO,RE]{\thepage}

\fancyhead[CO]{\scriptsize\MakeUppercase{\essaytitle}}
\fancyhead[CE]{\scriptsize\MakeUppercase{\essayauthor}}

\renewcommand{\headrulewidth}{0pt}
\thispagestyle{empty}

\newcommand{\pointsection}[1]{%
  \par\vspace{1.5em}%
  {\centering\small #1\par}%
  \vspace{1em}%
}

% Custom maketitle with slightly bold, uppercase title
\makeatletter
\renewcommand{\maketitle}{
  \begin{center}
    {\fontseries{bx}\normalsize\MakeUppercase\thetitle \par}
    {\footnotesize\MakeUppercase{by \theauthor}}
  \end{center}
}
\makeatother

\begin{document}

\maketitle
\vspace{1.5em}
\begin{center}
  \small\textit{Domine, veni et vide.}\\
  Lord, come and see.\\
  (John 11:34)
\end{center}
\vspace{1.5em}
\small
Shall I be permitted today to open a tomb in the presence of the Court? Will not these delicate eyes be offended by such a deadly sight? I do not think that Christians should refuse to attend this spectacle with Jesus Christ. For it was to him that the words of our Gospel were spoken: `Lord, come and see' where the body of Lazarus has been laid. It was he who commanded that the stone be rolled away. And it is he who seems to say to us in turn: come and see for yourselves. Jesus did not refuse to see that dead body, which was for him an object of pity and a subject for a miracle. Yet it is we, wretched mortals, who refuse to look upon this mournful scene, as the proof of our errors. Come and see with Jesus Christ, and let us forever be undeceived about all the goods that death strips away.

It is an odd weakness of mankind, that while death surrounds us in its myriad forms, it is never present to our minds. At funerals one only hears words of astonishment that a mortal man has died. Each brings to mind the last time he spoke with the deceased and what they had spoken about. Then, all of a sudden, he was dead. And we say: So this is what man is! But who is it that makes these observations? One who is himself a man; one who does not apply the lesson to himself; one who is not mindful of his own destiny. Or, if some transitory desire to prepare himself for death passes through the mind, he soon casts off such gloomy thoughts. It may even be said that mortals take no less care to bury the thoughts of death than they do the dead themselves. Yet perhaps these thoughts will have more of an effect in our hearts if we meditate upon them with Jesus Christ at the tomb of Lazarus. Let us ask him to imprint them upon our minds by the grace of his Holy Spirit, and let us strive to merit that grace through the intercession of the blessed Virgin. \textit{Ave Maria \ldots .}\footnote{Sermon for the Wednesday of the Fourth Week of Lent, preached in the royal chapel of the Louvre, in the presence of King Louis XIV, on March 22, 1662. The text is from Bossuet's \textit{Oeuvres Oratoires}, édition critique de l'abbé J. Lebarq, revue et augmentée par Ch. Urbain et E. Levesque (Paris: Desclée, 1926), IV: 262--81. Translation by Christopher O. Blum.}

Of all the passions of the human mind, one of the most violent is the desire to know. Our curiosity expends itself to find either some undiscovered secret in the order of nature, or some unknown skill in the works of art, or some unusual refinement in the conduct of our affairs. Yet in our keen desire to enrich our minds with new discoveries, we are like those who by looking far ahead fail to see the objects that surround them. By this I mean that our mind, devoting itself to great efforts upon things far afield and, as it were, wandering about the world, passes so quickly over things near at hand that we spend our entire lives not knowing the very things that touch upon us, and not only the things that touch upon us, but also what we are ourselves.

We have no more pressing need than to gather back to ourselves these straying thoughts. It is to this end, Christians, that I invite you today to accompany the Savior to the tomb of Lazarus: ``\textit{Veni et vide}; Come and see.'' O mortals, come and contemplate the spectacle of mortal things. O man, come to learn what you are.

You may be surprised that I would speak to you of death in order to teach you what you are, and you may think that man is misrepresented if he is depicted when he no longer exists. Yet if you carefully attend to what the tomb presents to us, you will surely agree that there is neither a more truthful interpreter, nor a more faithful mirror of humanity.

The nature of a composite being is never more distinctly to be perceived than in the dissolution of its parts. When joined, they mutually alter one another, and so they must be separated in order to be well understood. In fact, the fellowship of the body and the soul is such that the body seems to us to be more than it really is while the soul seems to be something less. But when, having been separated from one another, the body returns to the earth while the soul itself is put into a condition that allows it to return to heaven, whence it came, then we see each of them in its purity. We have, therefore, only to consider what death takes from us and what it leaves to us, which part of our being falls under its blows and which sustains itself amidst the ruin. Then we shall understand what man is, and so I shall not fear to affirm that it is in the bosom of death and in its dark shadow that an immortal light shines to illumine our minds about the condition of our nature. Let us run then, o mortals, and see what humanity is by gazing into the tomb of Lazarus. Come and see in one view the end of all your designs and the beginning of all your hopes. Come and see both the dissolution of your being and its renewal. Come and see the triumph of life amidst the victory of death: \textit{Veni et vide}.

O Death, we give you thanks for the light that you shine upon our ignorance. You alone can convince us of our lowliness; you alone can teach us of our dignity. If man thinks too well of himself, you know how to beat down his pride. If man disdains himself overmuch, you know how to rekindle his courage. And to bring all of these thoughts together in their proper balance, you teach him these two truths that open his eyes to accurate self-knowledge: that he is contemptible when he passes away, but infinitely worthy when he attains eternity. These two important considerations will form the subject of our discourse.

\pointsection{First Point}

It is a bold enterprise to tell man of his littleness. Each of us is jealous of what he is and would rather be blind than know his own weaknesses. Those who enjoy great wealth especially want to be treated delicately. They take no pleasure in the notice of their failings; they would prefer that, if we were to see them, at least we should keep them hidden from view. Nevertheless, thanks to death, we are at liberty to speak. From the vantage of death, none is so great in this world that he cannot recognize his littleness. Long live the Eternal One! O human greatness, from whatever side I look upon you--unless, of course, I consider you as having come from God and as owing everything to God, for in that case I discover in you a ray of Divinity that rightly attracts my regard. When, however, I look upon you as something purely human--I say it again, from whatever side I look upon you--I see nothing in you to esteem, because wherever I turn I always find myself facing death, which casts its shadow upon everything that worldly brilliance would illuminate, and then I no longer know to what I should attribute the august name of greatness, nor do I see anything worthy of being so designated.

Let us be convinced of this important truth, Christians, by an invincible course of reasoning. The accident cannot be nobler than the substance, nor the accessory more important than the principal, nor the building more secure than its foundation, nor, finally, can what is added to our being be greater or more important than our very being itself. Now, just what is our being? Tell us, O Death, for proud man no longer believes me. And yet you are mute, O Death; you speak only to our eyes. A great king will lend you his voice, so that our ears might hear you and our hearts might receive the most incisive of truths.

Here is the fine meditation that David entertained from the throne and amidst his court. It merits your attention, Sire: \textit{Ecce mensurabiles posuisti dies meos, et substantia mea tanquam nihilum ante te} (Ps. 38:6). O Eternal King of the ages! You are always with yourself and in yourself. Your being, forever enduring, neither passes away, nor changes, nor has limits; ``and behold, my days you have measured, and my substance is nothing before you.'' No, my substance is nothing before you, and every finite being is nothing, because what has limits has an end, and when it has arrived at that end, a final moment suffices to destroy everything as if it had never existed. What is a hundred years, what is a thousand, when they can be erased by a single moment? Multiply your days like the deer that fable or natural history would make to live for so many centuries. Live as long as the great oaks under which our ancestors lie and which will continue to shade our descendants. And fill that immense space with honors, riches, and pleasures. What shall any of it profit you, when death's final breath, weak and languid though it be, will cast down this vain display with the same ease as a house of cards, a mere child's plaything? What will it serve you to have written so much in this book, to have filled all its pages with beautiful script, when in the end a single stroke will erase it all? And does this one last stroke leave any trace even of itself? Or will not the last stroke of that final moment be itself lost in the great gulf of nothingness? There will not remain on earth even a single vestige of what we are: the flesh will change its nature, the body will take on another name, ``even that of cadaver will not long remain, for it will become,'' as Tertullian says, ``something that has no name in any language,'' so true it is that everything dies in it, including even those funereal terms with which we describe its wretched remains: \textit{Post totum ignobilitatis elogium, caducae in originem terram, et cadaveris nomen; et de isto quoque nominae periturae in nullum inde jam nomen, in omnis jam vocabuli mortem.}

What, then, is my substance, O great God? I come into life only to leave it soon after. I show myself like the others, but must disappear. Everything calls us unto death. Nature, almost envious of the gift she gives us, often tells us and shows us that she cannot allow us to keep for very long the little bit of matter she has loaned us, that it must not remain in the same hands, and that it must circulate forever. She has need of it for other forms; she requires it for other works.

The constant stream of recruits to mankind--that is, the children who are born--as they grow and advance, seem to push us from behind, saying, ``Get out of our way; now it is our turn.'' And just as we have seen others go before us, so others will see us pass, and they will present the same spectacle to their successors. O God! I ask again, what are we? If I look before me, what an infinite space in which I am not! And if I turn my gaze backward, what a fearful succession in which I am no longer! What a small place do I occupy in the immense abyss of time! I am nothing. So small an interval cannot distinguish me from nothingness. I was sent here only to make up the number; there was no need of me, and the play would have been no less well played had I remained behind the curtain.

Let us consider the matter again in more subtle terms. We see that it is not the length of our lives that distinguishes us from nothingness. And you know, Christians, that we are never separated from it by more than a single moment. Now we lay hold of one such moment, and now it dies. And with it, we should all die, if we did not immediately lay hold of another just like it. Until at last, there comes one which we cannot grasp, no matter what efforts we make to reach out for it. And then we shall immediately fall for lack of support. O fragile support of our being! O ruinous foundation of our substance! \textit{In imagine pertransit homo}. (Ps. 38:7) Truly, man passes like a shadow, or like an image in figure. And as he himself is nothing solid, so also does he pursue vain things, the image of the good, and not the good itself.

How little is the place that we occupy in this world! So little indeed and of such scant importance that it seems that my whole life is but a dream. Sometimes, with Arnobius, I wonder whether I sleep or wake: \textit{Vigilemus aliquando, an ipsum vigilare, quod dicitur, somni sit perpetui portio?} I do not know whether what I call waking is not perhaps a somewhat more excited part of a deep sleep, and whether I am seeing things that are real, or whether I am only troubled by fantasies and illusions. \textit{Praeterit figura huius mundi}. (I Cor. 7:31) The shape of this world is passing away, and, compared to God, my substance is nothing.

\pointsection{Second Point}

We must not doubt, Christians, that although we be relegated to this lowest part of the cosmos, to the theater of change and the empire of death, and although we carry it in our very bosom, nevertheless, through the shadowy knowledge that we gain from our prejudiced senses, if we know how to return to ourselves, we will find therein a principle whose vigor testifies to a heavenly origin, a principle that does not fear corruption.

I am not one of those who place great stock in human knowledge; and yet I confess that I cannot contemplate without admiration those marvelous discoveries by which science has sounded the depths of nature, nor the many fine inventions that art has found to make it suitable for our use. Man has almost changed the face of the earth. By his mind he has tamed animals that surpass him in strength. He has known how to discipline their brutish temper and channel their heedless liberty. He has even placed his mark upon the inanimate creatures. Has not his industry forced the earth to give him more suitable food, plants to adjust their wild bitterness in his favor, and even poisons to change into medicines for the love of him? It is needless to tell you how he has controlled the elements, after all of the miracles he daily performs with the most intractable of them, that is, with his two great enemies fire and water, who nevertheless agree to serve him in so many useful and necessary works. What else? He has climbed to the very heavens. In order to travel more safely, he has taught the stars to guide his voyages; to measure his life more accurately, he has obliged the sun to keep an account of all of his steps. Yet we must leave the long and scrupulous enumeration to rhetoric and be content as theologians to remark that God, intending man to be the head of the universe--as the oracle of Scripture says--formed him upon so noble a foundation that, although deformed by his crime, he has retained a certain instinct to seek what he lacks throughout the whole course of nature. That is why, if I may say so, he boldly delves everywhere, as if upon his own domain; and there is no part of the universe in which he has not set the sign of his industry.

How could such ascendancy have been gained by a creature with a body so weak and so liable to be assailed by every other one? It would not have happened unless his mind had a force superior to the whole of visible nature, an immortal breath of the Spirit of God, a ray of light from his face, an aspect of his countenance.

No, it could not be otherwise. If a skilled artisan makes a machine, no one is able to use it except with his instructions. God made the world as an immense machine that his wisdom alone could have invented and that his power alone could have built. O man, he established you to make use of it. He placed in your hands the whole of nature, that you might apply it to your ends. He even allowed you to decorate and embellish it by your art: for what is art but the embellishment of nature? You can add a little color to decorate this marvelous canvas; but how could you move even slightly a machine at once so strong and so delicate? Or in what way could you add even a single correct brush stroke to so rich a painting if you did not have in yourself and in some part of your being some art derived from that first art, some secondary ideas derived from the original ones, in a word, some resemblance, some outflow, some portion of that artisan Spirit that made the world? If this be true, Christians, who cannot see that all of nature drawn together is unable to extinguish so beautiful a ray of the power that sustains it, and that our soul, superior to the world and to all the powers that compose it, has nothing to fear except from its author?

Let us prolong this beneficial meditation upon the image of God in us, and let us see by what maxims that beloved creature man, destined to make use of all the others, prescribes for himself what he ought to do. I admit that, in our state of corruption, it is here that we discover our weakness. Nevertheless, I cannot fail to marvel at those unchanging moral rules that reason has laid down. What! This soul, immersed in the body, so closely married to all of the passions, who languishes, who is no longer herself when the body suffers: by what light has she been able to perceive that her felicity lies apart from it? By which she might boldly say--while all her senses, all her passions, and almost all of nature cry out against her--``Death, to me, is gain,'' (Philippians 1:21) and ``I rejoice in my afflictions.'' (Col. 1:24) Must it not be that she has discovered within herself a most exquisite beauty in what is called duty, in order to dare to affirm that for friends, country, prince, and altar one ought fearlessly to face immense labors, incredible sorrows, and the certitude of death? Is it not a kind of miracle that these maxims of courage, probity, and justice have not ever been abolished--I do not say by the passage of time, but by contrary practice--and that to the benefit of the human race there have always been many fewer persons who deny them than there have been those who have practiced them perfectly?

It may not be doubted that there is a divine light within us: ``A ray of light from your face, O Lord, has been lodged in our souls: \textit{Signatum est super nos lumen vultus tui, Domine}.'' (Ps. 4:7) Here we discover, as in a globe of light, an immortal charm in rectitude and virtue. It is the first Reason that reveals itself to us by its image. It is Truth himself who speaks to us and who causes us to understand that there is something in us that does not die, for God has made us capable of finding happiness even in death.

All of this is but nothing, Christians. For here is the most marvelous brush stroke of the divine likeness. God knows himself and contemplates himself. His life is to know himself, and because man is his image, he also wants man to know that He is eternal, immense, infinite, wholly immaterial, free from every limit, free from every imperfection. Christians, what is this miracle? We who sense only what has limits, who see nothing but what changes, whence have we been able to understand this eternity? Whence have we dreamt this infinity? ``O eternity, O infinity,'' says Saint Augustine, ``that our senses do not even suspect: how is it that you have entered our souls?'' But if all that we are is body and matter, how have we been able to conceive of pure spirit? And how have we been able even to invent the name?

I know what will be said at this point, and rightly: that, when we speak of these spirits we do not understand what we are saying. Our weak imagination, unable to sustain so pure an idea, always presents some little body with which to clothe it. Yet after reason has made its utmost effort to make those bodies subtle and fine, do you not sense that at the same time from the bottom of our soul there arises a heavenly light that clears away all those delicate phantoms by which we have imagined them? If you press onward and ask what it is, a voice rises from the center of the soul: I do not know what it is, yet it is certainly not that. What power, what energy, what secret strength does the soul find with which to correct itself, to undeceive itself, and to reject its own thoughts? Who does not see the hidden resource which has not yet used all of its strength, and which, though it be constrained, though it does not enjoy freedom of movement, nevertheless demonstrates by its vigor that it is not wholly bound to matter and that it is attached to some higher principle?

It is true, Christians, and I confess it: we cannot long maintain such noble ardor. The soul soon returns into her matter. The soul has her languor and frailty, and, as it were, her baseness, which--should she not be enlightened from elsewhere--almost force her to doubt what she is. This is why the worldly wise, in seeing man on the one hand so great and on the other so wretched, have not known what to think or to say: some making him out to be a god, others nothingness; some say that nature cherishes him like a mother and delights in him; others that, like an evil stepmother, she exposes him and considers him her castoff. And a third group, not knowing what to conclude from this mixture, responds that nature has been at play in uniting two parts that have nothing to do with one another, and so, by a kind of caprice, has formed this prodigy called man.

You will judge rightly, Christians, that none of these has hit the mark, and that only the Faith can solve so great a riddle. You mislead yourselves, O wise men of this age: man is not nature's delight, for she offends him in so many ways. Nor can man be her castoff, for there is something in him worth more than nature herself--I speak of the nature presented to us by the senses. Now, to speak of caprice in the works of God is to blaspheme against his wisdom. Yet whence so strange a disproportion? Must I tell you, Christians? Do not these ramshackle cabins with such magnificent foundations declare plainly enough that the work is not complete? Contemplate this grand edifice: you will see in it the marks of the divine hand. Yet the unevenness of the work will soon cause you to see that sin has mixed in its own. O God! What is this mixture? I am lost, and am almost ready to cry out with the prophet: ``\textit{Haeccine est urbs perfecti decoris, gaudium universae terrae?} Is this Jerusalem? Is this the city, is this the temple, the honor, the joy of the whole world?'' (Lamentations 2:15) And for my part I say: is this the man made in the image of God, the miracle of his wisdom, and the masterpiece of his hands?

Truly, it is he. Whence this discord? Why do I see such ill-assorted parts? Because man wanted to build upon his Creator's work in his own way, and he departed from the plan. And so, in opposition to the orderliness of the initial design, the immortal and the corruptible, the spiritual and the carnal, in a word, the angel and the beast found themselves joined together in one. Here is the clue to the riddle, here is the unraveling of the tangle: faith has restored us to ourselves, and our shameful failings can no longer hide our natural dignity.

But, alas, what does this dignity profit us? Even though we still breathe some air of greatness amidst our ruins, we are no less prostrate beneath them. Our former immortality serves only to make the tyranny of death all the more unbearable. And although our souls escape it, if they are made wretched by sin then they have nothing to boast about so burdensome an eternity. What shall we say then, Christians? How shall we reply to so pressing a complaint? Jesus Christ will respond to it in our Gospel. He comes to see the dead Lazarus; he comes to visit human nature groaning under the empire of death. This visit was not without cause: it is the artisan himself who comes in person to see what is lacking in the building, for he has a plan to renovate it in accord with his initial design: \textit{secundum imaginem ejus qui creavit illum}. (Col 3:10)

O soul, filled with crimes, rightly do you fear an immortality that would render your death eternal! Yet behold in the person of Jesus Christ the resurrection and the life. He who believes in him shall not die. He who believes in him is already living with a spiritual and interior life, living by the life of grace that draws after it the life of glory. But the body is nevertheless subject to death! O Soul, console yourself. If this divine architect who has undertaken to repair you allows the old structure of your body to fall piece by piece, it is so that he may return it to you in a better state and may rebuild it in a better order. For a little while, the body will fall under the empire of death, but it will not leave anything in his hands save mortality itself.

Do not let yourselves be convinced by medical reasoning that corruption is a natural result of composition and mixture. We must lift our minds higher and believe according to the principles of Christianity that what subjects the flesh to the necessity of corruption is that it is an attraction to evil, a source of evil desires, indeed, as the holy apostle says, a ``flesh of sin.'' (Romans 8:3) Such a flesh ought to be destroyed, I say even in the elect, because in the condition of a flesh of sin it does not deserve to be joined to the souls of the blessed nor to enter into the kingdom of God: \textit{caro et sanguis regnum Dei possidere non possunt}. (I Cor. 15:50) It must, therefore, change its initial form in order to be renewed and to lose its first condition in order to receive a second one from the hand of God. Like an old and haphazard building that is allowed to fall into disrepair in order to build a new one in a more lovely architectural order, so also this flesh, made unruly by sin and concupiscence, was allowed to fall into ruin in order that God may remake it in his way and according to the initial design of his creation. It must be reduced to dust because it gave service to sin.

Do you not see the divine Jesus who has the tomb opened? It is the Prince who opens the prison doors for the suffering captives. The dead bodies shut up within will one day hear his word, and they will be raised up like Lazarus. They will be raised up better than Lazarus, because they will be raised up never to die again, and because death, as the Holy Spirit says, will be drowned in the abyss, never to appear again: \textit{et mors ultra non erit}. (Apoc. 21:4)

What then do you fear, Christian soul, from death's approach? Perhaps in seeing your house fall you fear that you will lack shelter? But listen to the holy apostle: ``We know,'' we know, he says, we are not led to believe by uncertain guesses, but we know most assuredly and with complete certitude, ``that if this house of earth and mud in which we live is destroyed, we have another dwelling place prepared for us in heaven.'' (II Cor. 5:1) O merciful conduct of the one who anticipates our needs! He has a plan, as Saint John Chrysostom fittingly said, to repair the house he has given us. When he destroys it and casts it down in order to make it anew, we must move out. Yet he himself offers us his palace, and within it, gives us an apartment wherein we may await in peace the complete reconstruction of our former abode.
\end{document}
